\section{Actual entire tensor-power frames and their genuine transition units}

\begin{lemma}\label{lem:fc-actual-tensor-frame-unit-coefficient}
\lean{EnvelopeProperty.SchemeTensorPowerFrames.tensorFrame,EnvelopeProperty.SchemeTensorPowerFrames.unitSection_precompIso,EnvelopeProperty.SchemeTensorPowerFrames.tensorFrame_coefficient,EnvelopeProperty.SchemeTensorPowerFrames.tensorFrame_unitSection}
\leanok
\uses{lem:fc-original-tensor-section-homs,prop:generated-line-generic-frame-norms}
For actual entire original frames $e:P\simeq\mathcal O_X$ and
$d:Q\simeq\mathcal O_X$ on any original scheme, their actual
tensor frame is $e\otimes d$ followed by the true structure-unit
multiplication. Its coefficient on every genuine pure tensor
section is the product of the two original coefficients. Its
inverse image of one is exactly the genuine pure tensor of the
two original inverse-frame unit sections. Precomposing an actual
frame by an original module isomorphism transports its inverse
unit section by the actual inverse isomorphism.
\end{lemma}
\begin{proof}
Naturality of the actual original pure tensor section moves
$e\otimes d$ through that section, and genuine structure-unit
multiplication sends the resulting tensor to the product of
its original coefficients. Apply this to the two actual inverse
unit sections: both original coefficients equal one by the entire
frame inverse identities. Thus their pure tensor has coefficient
one in the actual tensor frame. Applying the inverse of that
whole frame identifies it with the inverse unit section. For
precomposition, invert the entire composite isomorphism; its
inverse is the original frame inverse followed by the inverse
of the preceding original isomorphism. Applying to one proves
the final assertion.
\end{proof}

\begin{definition}\label{def:fc-actual-entire-power-frame}
\lean{EnvelopeProperty.SchemeTensorPowerFrames.powerFrame,EnvelopeProperty.SchemeTensorPowerFrames.powerFrame_unitSection_succ}
\leanok
\uses{lem:fc-actual-tensor-frame-unit-coefficient,lem:bundle-tensor-power-isomorphism}
Every actual entire original frame $e:P\simeq\mathcal O_X$
induces an actual entire frame $e^{[n]}$ of the same explicit
original tensor power $P^{\otimes n}$ for every $n\geq0$.
The zeroth frame is the identity structure frame. Recursively,
the next frame uses the genuine explicit-to-monoidal tensor
comparison followed by the actual tensor frame of $e^{[n]}$
and $e$. Its inverse unit section is exactly the actual explicit
pure tensor of the preceding inverse unit section and that of $e$.
\end{definition}
\begin{proof}
The zeroth explicit tensor power is the actual structure module.
For a successor, the actual explicit tensor comparison identifies
the original power with the native tensor of the preceding power
and $P$. The preceding lemma constructs an entire original frame
of this native tensor. Composing these genuine isomorphisms gives
the asserted original power frame. The precomposition unit-section
identity and the genuine tensor-frame inverse-unit formula give
exactly the stated recursive original pure tensor section.
\end{proof}

\begin{theorem}\label{thm:fc-actual-power-frame-transition-units}
\lean{EnvelopeProperty.SchemeTensorPowerFrames.transitionUnit_precompIso,EnvelopeProperty.SchemeTensorPowerFrames.transitionUnit_tensorFrame,EnvelopeProperty.SchemeTensorPowerFrames.transitionUnit_powerFrame,EnvelopeProperty.SchemeTensorPowerFrames.transitionUnit_powerFrame_value}
\leanok
\uses{def:fc-actual-entire-power-frame,lem:fc-original-frame-transition-unit-value,lem:fc-original-frame-transition-unit-self}
On $\operatorname{Spec}A$ for any commutative ring $A$, precomposing
two actual entire frames by the same original module isomorphism
preserves their genuine transition unit. The transition between
two genuine tensor frames is the product of the transition units
of their two factors. Consequently, for any two actual original
frames $e_0,e_1$ of the same original module and every $n\geq0$,
\[
 u(e_0^{[n]},e_1^{[n]})=u(e_0,e_1)^n.
\]
Every actual valuation therefore gives
$\nu(u(e_0^{[n]},e_1^{[n]}))=\nu(u(e_0,e_1))^n$.
All frames and units are genuinely constructed; no power or
transition compatibility identity is supplied.
\end{theorem}
\begin{proof}
The transition unit is the second-frame coefficient of the first
frame's inverse unit section. For common precomposition, the
preceding original isomorphism and its inverse cancel on that
section, proving invariance of the actual unit. For tensor frames,
the preceding tensor inverse-unit and coefficient formulas express
this same actual coefficient as the product of the two transition
coefficients. The genuine affine section isomorphism preserves
multiplication; equality of values gives equality of the actual
units. Induct on $n$. At zero the two actual frames are identical,
so the transition is one. At a successor, common-precomposition
invariance removes the true explicit tensor comparison, and the
tensor-frame transition identity multiplies the preceding unit by
the degree-one unit. The induction hypothesis and the power
recurrence prove the asserted equality on genuine units. Apply
the actual multiplicative valuation and its preservation of powers
to obtain the final exact valuation formula.
\end{proof}
